\chapter{Advantage \& Disruption}\label{ch:advantage-and-disruption}

% Scope: resource-based view, VRIN, Blue Ocean, the innovator's dilemma.
% Cross-link: durable advantage is what sustains LTV > CAC over time (\cref{ch:customer-economics}).

Structural analysis (\cref{ch:competitive-analysis}) explains industry-level returns;
it cannot explain why some firms earn above-average returns \emph{inside} unattractive
industries for decades. The resource-based view, Blue Ocean Strategy, and the innovator's
dilemma each offer a different answer to that question --- and each implies a different
strategic response when the answer changes.

\section{The Resource-Based View \& VRIN}\label{sec:rbv-vrin}
% complexity: 3

Which of the firm's assets and capabilities will sustain a competitive advantage, and
which will be competed away? The audit must precede any growth bet. Industries compete for
customers; firms compete for resources. If the resources that drive performance can be
copied or bought, returns revert to the cost of capital regardless of market share. Durable
advantage requires resources that are, in a specific technical sense, not replicable by
rivals.

\textcite{barney1991resource} defines a resource as a source of sustained competitive
advantage only when it satisfies all four criteria simultaneously:
\begin{description}
  \item[Valuable] The resource enables the firm to exploit opportunities or neutralise threats
        in its environment --- i.e.\ it is linked to customer willingness to pay or cost
        reduction. A resource that passes no Five Forces test (\cref{ch:competitive-analysis})
        is not valuable in Barney's sense.
  \item[Rare] The resource is not possessed by a large number of current or potential
        competitors. Common resources generate competitive parity, not advantage.
  \item[Inimitable] Rivals cannot easily copy or acquire the resource. Inimitability stems
        from path dependence (accumulated via unrepeatable history), causal ambiguity
        (competitors cannot identify which resource causes performance), or social complexity
        (e.g.\ organisational culture, trust-based partnerships).
  \item[Non-substitutable] No strategically equivalent alternative resource exists. Even a
        rare, inimitable resource loses value if a substitute delivers the same strategic
        outcome.
\end{description}
These four criteria form a decision tree, not a checklist: fail any one and the resource
earns at most a transient advantage. Resources that pass all four generate
\emph{Ricardian rents} --- returns above the opportunity cost of the resource that persist
because the supply of the resource is inelastic.

VRIN is diagnostic, not prescriptive: it identifies where advantage currently lives but
offers no theory of how to build new resources. The framework also understates the
importance of \emph{dynamic capabilities} --- the firm's ability to reconfigure its resource
base as the environment shifts. A resource that is VRIN today may be rendered obsolete by
a technological discontinuity (\cref{sec:innovators-dilemma}), at which point Barney's
static audit returns the wrong answer.

\begin{example}
\emph{Decision: where should the SaaS firm invest to protect its margin?}

Audit of candidate resources against VRIN:

\begin{itemize}
  \item \textbf{Proprietary customer data (behaviour + integration logs)}: Valuable
        (enables personalisation and expansion revenue); Rare (each firm's data set is
        unique); Inimitable (path-dependent accumulation --- cannot be purchased wholesale);
        Non-substitutable (third-party intent data does not replicate longitudinal usage
        signals). \emph{Verdict: VRIN --- the defensible core.}
  \item \textbf{Integration ecosystem (50 + native connectors)}: Valuable and Rare at launch;
        Inimitable if integration depth exceeds a replication threshold; Non-substitutable
        if switching cost (\cref{ch:growth}) is high enough. \emph{Verdict: partially VRIN
        --- invest to deepen.}
  \item \textbf{Brand name}: Valuable (lowers CAC via organic search); not yet Rare in a
        fragmented market. \emph{Verdict: not VRIN --- supports, does not protect.}
  \item \textbf{Codebase}: Not rare (modern frameworks are standardised) and imitable with
        sufficient engineering spend. \emph{Verdict: competitive parity, not advantage.}
\end{itemize}

R\&D investment should deepen data infrastructure and integrations
(VRIN-passing resources) rather than rebuild generic features rivals already offer.
This directly raises LTV via lower churn (\cref{ch:customer-economics}) and is
the mechanism through which \cref{eq:roic} exceeds WACC sustainably.
\end{example}

\begin{admonition}{Warning}
Listing everything in the firm as a ``strategic resource.'' The VRIN criteria are
restrictive by design. Most resources are \emph{competitive necessities} --- required to
compete but conferring no advantage. Confusing a necessity (a CRM) with an advantage (the
proprietary data inside it) misdirects capital toward maintaining parity rather than
building rents.
\end{admonition}

\textcite{dranove7theconomics} synthesise RBV with industrial organisation: Ricardian rents
require both VRIN resources \emph{and} some barriers to resource market competition.
A resource cannot be VRIN in the product market if it is freely traded in the resource
market at its full NPV --- competitive resource markets dissipate rents before the product
market sees them.

The VRIN audit completes the structural picture that Five Forces began: \cref{ch:competitive-analysis}
identifies the threats from outside; VRIN identifies the resources that withstand them.
Sustainable ROIC exceeding WACC (\cref{eq:advantage-spread} below) is the financial
expression of a resource passing this test. \textcite{barney1991resource} is the
primary source; \textcite{dranove7theconomics} supply the IO integration.

\section{Blue Ocean Strategy}\label{sec:blue-ocean}
% complexity: 3

Rather than competing harder in an existing contested market, can we define a new value
curve that makes the competition temporarily irrelevant? Most strategy is about winning a
bigger share of an existing pie. Blue Ocean asks whether the pie can be reshaped by
simultaneously reducing cost \emph{and} raising buyer value --- eliminating features that
the industry competes on but customers do not value, and creating features that customers
want but no one offers.

\textcite{kim2015} propose the ERRC grid as the operational tool:
\begin{description}
  \item[Eliminate] Which factors the industry takes for granted should be removed entirely?
        These add cost without adding perceived value; removing them funds the ``Create''
        column.
  \item[Reduce] Which factors should be reduced well below the industry standard?
        Competitors over-deliver on these; the cost saving is re-deployable.
  \item[Raise] Which factors should be raised well above the industry standard?
        These address genuine unmet needs.
  \item[Create] Which factors should be introduced that the industry has never offered?
        This is the demand-creation mechanism; it targets non-consumption rather than
        poaching from rivals.
\end{description}
The result is a \emph{value curve} --- a plot of offering level across competitive factors
--- that diverges visibly from the industry average. Divergence is the test: a curve that
merely scores higher everywhere is not a Blue Ocean move, it is a premium strategy.

The ERRC grid is directional and qualitative; it provides no mechanism for estimating
whether the new value curve generates sufficient demand to cover the investment. A redesigned
offer that no customer segment values passes the ERRC test but fails commercially.
Pair with the segmentation lens of \cref{ch:segmentation-targeting-positioning} to verify
which non-consumers the move actually reaches.

\begin{example}
\emph{Decision: can the SaaS firm create a Blue Ocean relative to Salesforce and similar
enterprise CRMs?}

\begin{itemize}
  \item \textbf{Eliminate}: Complex implementation (6--12 month deployments), dedicated
        admin roles, on-premise data options.
  \item \textbf{Reduce}: Price relative to enterprise CRM (Salesforce runs
        \money{150}--\money{300}/user/month; our \money{50}/seat is a deliberate undercut),
        feature set relative to the full Salesforce Sales Cloud.
  \item \textbf{Raise}: Time-to-first-value (target: live in one day), native integrations
        with modern productivity tools, transparent usage-based billing.
  \item \textbf{Create}: Self-serve onboarding with no sales touchpoint required; in-product
        guided workflows built for non-technical operators.
\end{itemize}

The curve targets \emph{non-consumers} of CRM --- SMBs that could not absorb enterprise
implementation cost --- rather than Salesforce's existing customers. If successful, CAC
from this segment is lower (no sales team required for self-serve) and churn is lower
(faster time-to-value). Both improve the LTV:CAC ratio in \cref{ch:customer-economics}.
\end{example}

\begin{admonition}{Warning}
Applying ERRC post-hoc to rationalise existing positioning rather than as a \emph{forward-
looking design tool}. Critics (\parencite{dranove7theconomics}) note that most Blue Ocean
case studies are written after the fact, when the value curve that happened to succeed gets
labelled a deliberate ERRC move. Used rigorously the grid is a hypothesis about demand;
validate with market experiments before committing capital.
\end{admonition}

\begin{admonition}{Contested}
Blue Ocean is the most contested strategic framework in this book. The critics' core
objection: successful Blue Ocean cases are distinguishable from lucky product bets only
in retrospect. The ERRC grid does not tell you \emph{which} factors to eliminate or
create --- that requires demand-side insight the framework does not supply. Proponents
respond that the value-curve visualisation disciplines a management team away from
incremental thinking. Both sides are right; treat Blue Ocean as a creativity prompt
validated by customer evidence, not as a planning methodology.
\end{admonition}

Blue Ocean Strategy intersects with the jobs-to-be-done (\cref{ch:customer-economics})
view of demand: the ``Create'' column succeeds when it addresses a job no existing product
does well. The mechanism is closer to disruptive innovation (below) than to traditional
differentiation, because it targets different success criteria rather than outperforming
incumbents on shared ones.

The ERRC grid answers the design question; the VRIN audit (\cref{sec:rbv-vrin}) then asks
whether the resulting offer is defensible once rivals can observe it. A Blue Ocean that is
VRIN-passing is durable; one that is not will be imitated and the advantage compressed.
\textcite{kim2015} is the primary source.

\section{The Innovator's Dilemma}\label{sec:innovators-dilemma}
% complexity: 5 -> figure disruption-scurve

Is an apparently inferior entrant a genuine disruption threat, or merely a different product?
The answer determines whether to respond aggressively, ignore, or acquire. Incumbents lose
not because they are badly managed but because they are well managed: they listen to their
best customers, invest where margins are highest, and improve performance on the dimensions
their customers most care about. Disruptors enter on dimensions no current customer asks for,
serve non-consumers or the underserved, and improve until they overtake the mainstream. The
danger is invisible until it is too late.

\textcite{christensen1997} identifies two innovation types:
\begin{description}
  \item[Sustaining innovations] Improve performance on dimensions existing customers value.
        Incumbents almost always win these contests; they have scale, established channels,
        and the right incentives.
  \item[Disruptive innovations] Introduce a different performance dimension --- often
        initially lower on the established metric (performance, features, reliability) but
        superior on accessibility, price, or simplicity. They enter in non-consumption or
        low-end segments the incumbent does not defend.
\end{description}
The disruption mechanism is an S-curve mismatch: the incumbent's sustaining trajectory
overshoots the mainstream market's demand ceiling. The disruptor's improvement trajectory
starts below the threshold but, as it rises on the S-curve of its own technology, it
eventually crosses the mainstream performance threshold --- at which point the incumbent's
customer base switches en masse because the disruptor is now ``good enough'' and cheaper.

\begin{figure}[htbp]
  \centering
  \includegraphics[width=0.86\linewidth]{disruption-scurve}
  \caption{The innovator's dilemma as dual S-curves. The incumbent's trajectory overshoots
    the mainstream demand threshold; the disruptor enters below it in an underserved segment
    and crosses the threshold as its S-curve matures. The crossing point is when disruption
    occurs.}
  \label{fig:disruption-scurve}
\end{figure}

Christensen's criterion for disruption is entry in \emph{underserved or non-consumption}
segments --- not just ``entry from below'' in a price sense. Every incumbent calls every
low-price entrant ``disruptive''; the framework's predictive power depends on correctly
identifying whether the entrant's performance trajectory will eventually satisfy mainstream
demand. Many low-end entrants remain confined to low-end segments permanently (low-end
encroachment, not disruption). The test is trajectory, not current price.

\begin{example}
\emph{Decision: should the SaaS firm treat no-code automation tools (Zapier, Make,
Airtable) as a disruption threat?}

No-code tools currently score low on the dimensions the SaaS's existing customers value
(data depth, reporting, compliance). They are not competitive \emph{today} for the
mid-market buyer.

The disruptive question is: is the no-code S-curve improving fast enough to reach
mid-market performance requirements within a planning horizon? Evidence suggests yes:
Airtable's 2022 funding explicitly targets mid-market workflows. Christensen's test is
whether no-code tools are improving on the \emph{right} trajectory, not whether they are
competitive today.

Response: invest in the VRIN resources (data depth, integrations) that no-code tools
cannot replicate quickly, while watching the no-code improvement rate as a leading indicator.
A two-year lag between ``enters non-consumption'' and ``crosses mainstream threshold'' is
historically common --- that is the response window.
\end{example}

\begin{admonition}{Warning}
Every incumbent labels every entrant ``disruptive'' because the word has become a compliment.
\emph{Disruptive} in Christensen's sense is a precise technical term: it means the entrant
initially competes on different performance criteria in a different market segment. A
premium entrant that targets the incumbent's best customers is a \emph{sustaining} threat
(harder to beat on features, not disruption). Misclassifying the threat type produces the
wrong response: disruption requires a separate low-end business unit; sustaining threats
require direct product investment.
\end{admonition}

The dilemma is organisational as well as technological: a rational incumbent allocates
resources toward its highest-margin customers, precisely the ones the disruptor does not
initially target. Christensen's solution --- a structurally separate unit with different
cost structure and metrics --- is a corporate-governance problem (\cref{ch:organizational-behavior}).
The financial expression of the dilemma is an ROIC that looks excellent right up until the
disruption crossing, then collapses (\cref{eq:advantage-spread}).

Disruption resets the VRIN audit: a previously inimitable resource (physical distribution,
in the case of music) can be substituted by a new technological infrastructure. The chapter's
final section ties the disruption dynamic to the financial metric that makes it legible.
\textcite{christensen1997} is the primary source.

\section{Advantage as a Value Spread}\label{sec:advantage-spread}
% complexity: 3

How do we know whether the firm currently possesses a competitive advantage, and whether
it is widening or narrowing over time? Qualitative frameworks diagnose the \emph{source} of
advantage; a financial metric is needed to confirm its \emph{magnitude} and track whether
strategy is working. The natural measure is whether returns exceed the cost of capital by
enough to justify the risk --- the same test a shareholder would apply.

Competitive advantage is operationally defined as a sustained positive value spread:
\begin{equation}\label{eq:advantage-spread}
  \text{Advantage spread} \;=\; \mathrm{ROIC} - \mathrm{WACC},
\end{equation}
where ROIC is defined in \cref{eq:roic} and WACC in \cref{eq:wacc}. When the spread is
positive, the firm earns Ricardian rents and creates economic profit (\cref{eq:economic-profit}).
When it is zero or negative, the firm earns competitive parity or destroys value regardless
of accounting profit.

The spread connects directly to customer economics: sustained ROIC $>$ WACC requires
sustained LTV $>$ CAC (\cref{eq:ltv-cac} in \cref{ch:customer-economics}), because the
net present value of every cohort must be positive and growing. A widening advantage spread
is the aggregate expression of a deepening moat.

ROIC can be manipulated through accounting choices (expensing vs capitalising R\&D, off-balance-sheet
liabilities). A firm that capitalises software development inflates the invested capital base
and suppresses ROIC; one that expenses it deflates the base. Normalise before comparing
across firms or across time. Similarly, a single-year positive spread may reflect a cyclical
peak rather than structural advantage --- the spread must be persistent to be evidence of a
moat.

\begin{example}
\emph{Decision: does the SaaS firm currently have a measurable advantage?}

From the fact sheet: NOPAT $\approx$\money{600000}, invested capital $\approx$\money{3000000},
giving ROIC $= \money{600000}/\money{3000000} = \qty{20}{\percent}$. WACC $= \qty{11}{\percent}$
(\cref{eq:wacc}).
\[
  \text{Advantage spread} = \qty{20}{\percent} - \qty{11}{\percent} = \qty{9}{\percent}.
\]
The spread is positive and material. But it must be assessed against the moat analysis:
if the VRIN audit shows only transient advantages, the spread will compress as rivals
replicate. A \qty{9}{\percent} spread sustained for five years at \money{3000000}
invested capital creates $5 \times 0.09 \times \money{3000000} = \money{1350000}$
of cumulative economic profit --- the reward for genuine VRIN resources.

If moat investment (\cref{ch:growth}) raises LTV by \money{500} (e.g.\ lower churn from
deeper integrations), the firm can profitably spend \money{500} more per customer
acquired without violating the LTV:CAC floor (\cref{eq:ltv-cac}).
\end{example}

\begin{admonition}{Warning}
Comparing ROIC to the accounting cost of debt rather than to WACC. Debt is cheap in nominal
terms; WACC blends in the much higher required return on equity. A firm earning
\qty{8}{\percent} ROIC that funds itself at \qty{5}{\percent} nominal debt \emph{appears}
profitable but destroys equity value if WACC $=\qty{11}{\percent}$.
\end{admonition}

McKinsey's valuation framework (\parencite{dranove7theconomics}) shows that market value
is the discounted present value of future economic profits, not accounting profits.
A firm trading at a premium to book value has positive NPV of future advantage spread;
a firm trading at a discount signals the market's expectation that ROIC will fall below WACC.
The EV multiples in \cref{ch:valuation} are the market's forecast of future spread, not
a snapshot of current spread.

The advantage spread (\cref{eq:advantage-spread}) ties the qualitative frameworks of this
chapter --- VRIN, Blue Ocean, disruption --- to the financial language of
\cref{ch:corporate-finance-core} and \cref{ch:valuation}. Sustained spread is the
empirical test that distinguishes durable advantage from temporary good fortune.
\cref{ch:growth} then asks how to widen it through moat-building.
